% Options for packages loaded elsewhere
\PassOptionsToPackage{unicode}{hyperref}
\PassOptionsToPackage{hyphens}{url}
\PassOptionsToPackage{dvipsnames,svgnames,x11names}{xcolor}
%
\documentclass[
<<<<<<< HEAD:inst/rmarkdown/templates/z-article/skeleton/skeleton.tex
11pt
twocolumn
]{article}

\usepackage[]{mathpazo}

\usepackage{amssymb,amsmath}
\usepackage{ifxetex,ifluatex}
\usepackage{fixltx2e} % provides \textsubscript
\ifnum 0\ifxetex 1\fi\ifluatex 1\fi=0 % if pdftex
  \usepackage[T1]{fontenc}
  \usepackage[utf8]{inputenc}
  \usepackage{textcomp} % provide euro and other symbols
\else % if luatex or xelatex
  \ifxetex
    \usepackage{mathspec}
    \usepackage{xltxtra, xunicode}
  \else
    \usepackage{fontspec}
  \fi
  \defaultfontfeatures{Mapping=tex-text, Scale=MatchLowercase}
  \newcommand{\euro}{€}

\fi



% use upquote if available, for straight quotes in verbatim environments
\IfFileExists{upquote.sty}{\usepackage{upquote}}{}
% use microtype if available
\IfFileExists{microtype.sty}{%
\usepackage{microtype}
\UseMicrotypeSet[protrusion]{basicmath} % disable protrusion for tt fonts
}{}
\usepackage[hmarginratio=1:1, top=32mm, columnsep=20pt, includeheadfoot]{geometry}

=======
  11pt,
  twocolumn]{article}
\usepackage{amsmath,amssymb}
\usepackage{iftex}
\ifPDFTeX
  \usepackage[T1]{fontenc}
  \usepackage[utf8]{inputenc}
  \usepackage{textcomp} % provide euro and other symbols
\else % if luatex or xetex
  \usepackage{unicode-math} % this also loads fontspec
  \defaultfontfeatures{Scale=MatchLowercase}
  \defaultfontfeatures[\rmfamily]{Ligatures=TeX,Scale=1}
\fi
\usepackage{lmodern}
\ifPDFTeX\else
  % xetex/luatex font selection
  \setmainfont[]{Times New Roman}
\fi
% Use upquote if available, for straight quotes in verbatim environments
\IfFileExists{upquote.sty}{\usepackage{upquote}}{}
\IfFileExists{microtype.sty}{% use microtype if available
  \usepackage[]{microtype}
  \UseMicrotypeSet[protrusion]{basicmath} % disable protrusion for tt fonts
}{}
>>>>>>> 68f0abee1e0db71c1f585da2585915c8eb49dbf6:inst/rmarkdown/templates/zarticle/skeleton/skeleton.tex
\makeatletter
\@ifundefined{KOMAClassName}{% if non-KOMA class
  \IfFileExists{parskip.sty}{%
    \usepackage{parskip}
  }{% else
    \setlength{\parindent}{0pt}
    \setlength{\parskip}{6pt plus 2pt minus 1pt}}
}{% if KOMA class
  \KOMAoptions{parskip=half}}
\makeatother
<<<<<<< HEAD:inst/rmarkdown/templates/z-article/skeleton/skeleton.tex


\usepackage{xcolor}

\makeatletter
\@ifpackageloaded{hyperref}{}{%
\ifxetex
  \usepackage[setpagesize=false, % page size defined by xetex
              unicode=false, % unicode breaks when used with xetex
              xetex]{hyperref}
\else
  \usepackage[unicode=true]{hyperref}
\fi
}
\@ifpackageloaded{color}{
    \PassOptionsToPackage{usenames,dvipsnames}{color}
}{%
    \usepackage[usenames,dvipsnames]{color}
}
\makeatother
\hypersetup{breaklinks=true,
            bookmarks=true,
            pdfauthor={Author Name () and  (University of Someplace)},
            pdftitle={Article Title: Subtitle},
            colorlinks=true,
            citecolor=blue,
            urlcolor=blue,
            linkcolor=blue,
            pdfborder={0 0 0}}
\urlstyle{same}  % don't use monospace font for urls
\usepackage[]{biblatex}
\addbibresource{sample.bib}







% Set default figure placement to htbp
\makeatletter
\def\fps@figure{htbp}
\makeatother



\setlength{\emergencystretch}{3em} % prevent overfull lines
\providecommand{\tightlist}{%
  \setlength{\itemsep}{0pt}\setlength{\parskip}{0pt}}



\setlength{\parindent}{0pt}
\setlength{\parskip}{6pt plus 2pt minus 1pt}
\setlength{\emergencystretch}{3em}  % prevent overfull lines
\setcounter{secnumdepth}{-\maxdimen} % remove section numbering



%%% Use protect on footnotes to avoid problems with footnotes in titles
\let\rmarkdownfootnote\footnote%
\def\footnote{\protect\rmarkdownfootnote}

\IfFileExists{bookmark.sty}{\usepackage{bookmark}}{\usepackage{hyperref}}
\IfFileExists{xurl.sty}{\usepackage{xurl}}{} % add URL line breaks if available
\urlstyle{same} % disable monospaced font for URLs

=======
\usepackage{xcolor}
\setlength{\emergencystretch}{3em} % prevent overfull lines
\providecommand{\tightlist}{%
  \setlength{\itemsep}{0pt}\setlength{\parskip}{0pt}}
\setcounter{secnumdepth}{-\maxdimen} % remove section numbering
\newlength{\cslhangindent}
\setlength{\cslhangindent}{1.5em}
\newlength{\csllabelwidth}
\setlength{\csllabelwidth}{3em}
\newlength{\cslentryspacingunit} % times entry-spacing
\setlength{\cslentryspacingunit}{\parskip}
\newenvironment{CSLReferences}[2] % #1 hanging-ident, #2 entry spacing
 {% don't indent paragraphs
  \setlength{\parindent}{0pt}
  % turn on hanging indent if param 1 is 1
  \ifodd #1
  \let\oldpar\par
  \def\par{\hangindent=\cslhangindent\oldpar}
  \fi
  % set entry spacing
  \setlength{\parskip}{#2\cslentryspacingunit}
 }%
 {}
\usepackage{calc}
\newcommand{\CSLBlock}[1]{#1\hfill\break}
\newcommand{\CSLLeftMargin}[1]{\parbox[t]{\csllabelwidth}{#1}}
\newcommand{\CSLRightInline}[1]{\parbox[t]{\linewidth - \csllabelwidth}{#1}\break}
\newcommand{\CSLIndent}[1]{\hspace{\cslhangindent}#1}
\ifLuaTeX
  \usepackage{selnolig}  % disable illegal ligatures
\fi
\IfFileExists{bookmark.sty}{\usepackage{bookmark}}{\usepackage{hyperref}}
\IfFileExists{xurl.sty}{\usepackage{xurl}}{} % add URL line breaks if available
\urlstyle{same}
>>>>>>> 68f0abee1e0db71c1f585da2585915c8eb49dbf6:inst/rmarkdown/templates/zarticle/skeleton/skeleton.tex
\hypersetup{
  pdftitle={Article Title},
  colorlinks=true,
  linkcolor={blue},
  filecolor={Maroon},
  citecolor={blue},
  urlcolor={blue},
  pdfcreator={LaTeX via pandoc}}

<<<<<<< HEAD:inst/rmarkdown/templates/z-article/skeleton/skeleton.tex






%% --------- above is almost identical with default rmarkdown-------------------
%-------------------------------------------------------------------------------
% TITLE SECTION
%-------------------------------------------------------------------------------

%%% Change title format to be more compact
\usepackage{titling}


\setlength{\droptitle}{-5\baselineskip} % Move the title up
% alternatively
% \setlength{\droptitle}{-2em}

% Set up for small logo in top corner of document
% alternative is flushright
\pretitle{
\begin{flushleft}
\includegraphics[height=2.2cm, keepaspectratio]{umd_logo.png}
\end{flushleft}
\setlength{\droptitle}{-5\baselineskip} % Move the title up
\begin{center}\huge\bfseries % Article title formatting
}

\posttitle{\end{center}} % Article title closing formatting

\title{Article Title
  % \pretitle{\vspace{\droptitle}}
\usepackage{etoolbox}
\makeatletter
\providecommand{\subtitle}[1]{% add subtitle to \maketitle
  \apptocmd{\@title}{\par {\large #1 \par}}{}{}}
\makeatother
\subtitle{Subtitle}\thanks{Thanks to all the important people.
Replication files are available on request.}


  \author{Author Name () \and

  
   (University of Someplace)

  \thanks{Thanks to all the important people. Replication files are
available on request.}


  \email{\href{\href{mailto:your.email@address.edu}{\nolinkurl{your.email@address.edu}}}{\href{mailto:your.email@address.edu}{\nolinkurl{your.email@address.edu}}} \\[1ex]
  % \preauthor{\centering\large\emph}
  % \postauthor{\par}
  % \preauthor{}\postauthor{}

  \predate{\par\begin{center}\large\vskip -1em}
  \postauthor{\end{center}}
  \date{December 15, 2022}
  % \predate{}\postdate{}



% add an option for endnotes ---

\usepackage{endnotes}
\renewcommand{\enotesize}{\normalsize}
\let\footnote=\endnote

\makeatletter
\def\@seccntformat#1{\llap{\csname the#1\endcsname \hspace{.075 in}}}
\makeatother

% % suppress chapter in section numbering
% \renewcommand*\thesection{\arabic{section}}

% spacing of bulleted lists
\usepackage{enumitem} %control spacing of enumerated items
\setlist[itemize]{noitemsep, topsep=0pt}

\usepackage{abstract} % Allows abstract customization
% Set the "Abstract" text to bold
\renewcommand{\abstractnamefont}{\normalfont\bfseries} % Set the "Abstract" text to bold
\renewcommand{\abstracttextfont}{\normalfont\small\itshape} % Set the abstract itself to small italic text


\newlength\tbspace
\setlength\tbspace{.25in}

% units.
%\usepackage{units}


\def\tightlist{}


%-------------------------------------------------------------------------------
%% document formatting
%-------------------------------------------------------------------------------


 % colors for tables and text
\usepackage{ragged2e} % justifying text
\usepackage{setspace} % spacing commands, automatically makes captions single-spaced
%  \setstretch{1.2} \frenchspacing
\usepackage{lastpage} % access number of last page for numbering in margin

%% fonts setup
\usepackage{gensymb} %degree symbol
\usepackage{array}
\usepackage{multirow}
\usepackage{wrapfig}
\usepackage{float}
\usepackage{colortbl}
\usepackage{pdflscape}
\usepackage{tabu}
\usepackage{threeparttable}
\usepackage{threeparttablex}
\usepackage{makecell}
\usepackage[hang]{footmisc}
\usepackage{booktabs}



\usepackage{titlesec} % Allows customization of titles
% removes numbered sections
\titleformat*{\section}{\sc \bf \large \color[rgb]{0,0,0}}
\titleformat*{\subsection}{\bf \large}
\titleformat*{\subsubsection}{\bf \normalsize}
\titlespacing\section{0pt}{8pt plus 2pt minus 2pt}{0pt plus 2pt minus 2pt}
\titlespacing\subsection{0pt}{6pt plus 2pt minus 2pt}{-2pt plus 2pt minus 2pt}
\titlespacing\subsubsection{0pt}{2pt plus 2pt minus 2pt}{-3pt plus 2pt minus 2pt}

\robustify\tnote

% table/figure captions (titles and legends)

% \usepackage{caption}
% % otherwise just
% \captionsetup{font=small, labelfont={sc, bf}}
% 
% table/figure captions (titles and legends)
\usepackage[singlelinecheck=false]{caption} % left justify table captions
\captionsetup{font={color=black, sf}, skip=4pt, size=small, labelfont=sf, labelsep = period} % color and spacing between caption and table proper, updated

\makeatletter
\let\runtitle\@title
\makeatother


=======
\title{Article Title\thanks{Thanks to all the important people.
Replication files are available on request.}}
\usepackage{etoolbox}
\makeatletter
\providecommand{\subtitle}[1]{% add subtitle to \maketitle
  \apptocmd{\@title}{\par {\large #1 \par}}{}{}
}
\makeatother
\subtitle{Subtitle}
\author{true \and true}
\date{February 03, 2023}

%%%%%%%%%%%%%%%%%%%%%%%%%

%\usepackage{blindtext} % Package to generate dummy text throughout this template

\usepackage[sc]{mathpazo} % Use the Palatino font
\usepackage[T1]{fontenc} % Use 8-bit encoding that has 256 glyphs
\linespread{1.05} % Line spacing - Palatino needs more space between lines
\usepackage{microtype} % Slightly tweak font spacing for aesthetics

\usepackage[english]{babel} % Language hyphenation and typographical rules

\usepackage[hmarginratio=1:1,top=32mm,columnsep=20pt]{geometry} % Document margins
\usepackage[hang, small,labelfont=bf,up,textfont=it,up]{caption} % Custom captions under/above floats in tables or figures
\usepackage{booktabs} % Horizontal rules in tables

\usepackage{lettrine} % The lettrine is the first enlarged letter at the beginning of the text

\usepackage{enumitem} % Customized lists
\setlist[itemize]{noitemsep} % Make itemize lists more compact

\usepackage{abstract} % Allows abstract customization
\renewcommand{\abstractnamefont}{\normalfont\bfseries} % Set the "Abstract" text to bold
\renewcommand{\abstracttextfont}{\normalfont\small\itshape} % Set the abstract itself to small italic text

\usepackage{titlesec} % Allows customization of titles
\renewcommand\thesection{\Roman{section}} % Roman numerals for the sections
\renewcommand\thesubsection{\roman{subsection}} % roman numerals for subsections
\titleformat{\section}[block]{\large\scshape\centering}{\thesection.}{1em}{} % Change the look of the section titles
\titleformat{\subsection}[block]{\large}{\thesubsection.}{1em}{} % Change the look of the section titles
>>>>>>> 68f0abee1e0db71c1f585da2585915c8eb49dbf6:inst/rmarkdown/templates/zarticle/skeleton/skeleton.tex

\usepackage{fancyhdr} % Headers and footers
\pagestyle{fancy} % All pages have headers and footers
\fancyhead{} % Blank out the default header
\fancyfoot{} % Blank out the default footer
<<<<<<< HEAD:inst/rmarkdown/templates/z-article/skeleton/skeleton.tex
\fancyhead[R]{\author{truetrue}} % Custom header text
\fancyfoot[RO,LE]{\thepage} % Custom footer text

% \fancypagestyle{firststyle}
% {
%    \fancyhf{}
%    %    \renewcommand{\footrulewidth}{0pt}
%    % }

\setlength\parindent{0pt}


% Redefines (sub)paragraphs to behave more like sections
\ifx\paragraph\undefined\else
\let\oldparagraph\paragraph
\renewcommand{\paragraph}[1]{\oldparagraph{#1}\mbox{}}
\fi
\ifx\subparagraph\undefined\else
\let\oldsubparagraph\subparagraph
\renewcommand{\subparagraph}[1]{\oldsubparagraph{#1}\mbox{}}
\fi


\usepackage{tikz}

\newcommand{\shrug}[1][]{%
\begin{tikzpicture}[baseline,x=0.8\ht\strutbox,y=0.8\ht\strutbox,line width=0.125ex,#1]
\def\arm{(-2.5,0.95) to (-2,0.95) (-1.9,1) to (-1.5,0) (-1.35,0) to (-0.8,0)};
\draw \arm;
\draw[xscale=-1] \arm;
\def\headpart{(0.6,0) arc[start angle=-40, end angle=40,x radius=0.6,y radius=0.8]};
\draw \headpart;
\draw[xscale=-1] \headpart;
\def\eye{(-0.075,0.15) .. controls (0.02,0) .. (0.075,-0.15)};
\draw[shift={(-0.3,0.8)}] \eye;
\draw[shift={(0,0.85)}] \eye;
% draw mouth
\draw (-0.1,0.2) to [out=15,in=-100] (0.4,0.95);
\end{tikzpicture}}




\linespread{1.05}

%-------------------------------------------------------------------------------
% Begin Document
%-------------------------------------------------------------------------------

\begin{document}

\maketitle


=======
\fancyhead[C]{Article Title $\bullet$ February 03,
2023 $\bullet$ truetrue} % Custom header text
\fancyfoot[R]{\thepage} % Custom footer text

\usepackage{titling} % Customizing the title section

\usepackage[backend=biber,style=authoryear,dashed=false,uniquename=false]{biblatex}
 %using biblatex, gotta run pdfLaTeX, then biber, then pdfLaTeX again
\addbibresource{sample.bib} %or whatever .bib is being used



\ifluatex
  \usepackage{selnolig}  % disable illegal ligatures
\fi


\usepackage{hyperref} % For hyperlinks in the PDF

\makeatletter
\@ifpackageloaded{hyperref}{}{%
\ifxetex
  \PassOptionsToPackage{hyphens}{url}\usepackage[setpagesize=false, % page size defined by xetex
              unicode=false, % unicode breaks when used with xetex
              xetex]{hyperref}
\else
  \PassOptionsToPackage{hyphens}{url}\usepackage[draft,unicode=true]{hyperref}
\fi
}

\@ifpackageloaded{color}{
    \PassOptionsToPackage{usenames,dvipsnames}{color}
}{%
    \usepackage[usenames,dvipsnames]{color}
}

\makeatother
\hypersetup{breaklinks=true,
            bookmarks=true,
            pdfauthor={Author Name () and  (University of Someplace)},
             pdfkeywords = {},
            pdftitle={Article Title: Subtitle},
            colorlinks=true,
            citecolor=blue,
            urlcolor=blue,
            linkcolor=blue,
            pdfborder={0 0 0}}
\urlstyle{same}  % don't use monospace font for urls

\setlength{\emergencystretch}{3em} % prevent overfull lines
\providecommand{\tightlist}{%
  \setlength{\itemsep}{0pt}\setlength{\parskip}{0pt}}
\setcounter{secnumdepth}{-\maxdimen} % remove section numbering


\usepackage{endnotes}
\renewcommand{\enotesize}{\normalsize}
\let\footnote=\endnote



%----------------------------------------------------------------------------------------
%	TITLE SECTION
%----------------------------------------------------------------------------------------
\setlength{\droptitle}{-5\baselineskip} % Move the title up

\pretitle{\begin{center}\Huge\bfseries} % Article title formatting
\posttitle{\end{center}} % Article title closing formatting
\title{Article Title\thanks{Thanks to all the important people.
Replication files are available on request.}}
\providecommand{\subtitle}[1]{}
\subtitle{Subtitle}
 % Article title
\author{%
\textsc{}\thanks{Thanks to all the important people. Replication files
are available on request.} \\[1ex] % Your name
\normalsize  \\ % Your institution
\normalsize \href{mailto:\href{mailto:your.email@address.edu}{\nolinkurl{your.email@address.edu}}}{\href{mailto:your.email@address.edu}{\nolinkurl{your.email@address.edu}}} % Your email address
%\and % Uncomment if 2 authors are required, duplicate these 4 lines if more
%\textsc{Jane Smith}\thanks{Corresponding author} \\[1ex] % Second author's name
%\normalsize University of Utah \\ % Second author's institution
%\normalsize \href{mailto:jane@smith.com}{jane@smith.com} % Second author's email address
}
\date{February 03, 2023} % Leave empty to omit a date
\renewcommand{\maketitlehookd}{%
>>>>>>> 68f0abee1e0db71c1f585da2585915c8eb49dbf6:inst/rmarkdown/templates/zarticle/skeleton/skeleton.tex
\begin{abstract}
Lorem ipsum dolor sit amet, consectetur adipiscing elit. Donec sit amet
libero justo. Pellentesque eget nibh ex. Aliquam tincidunt egestas
lectus id ullamcorper. Proin tellus orci, posuere sed cursus at,
bibendum ac odio. Nam consequat non ante eget aliquam. Nulla facilisis
tincidunt elit. Nunc hendrerit pellentesque quam, eu imperdiet ipsum
porttitor ut. Interdum et malesuada fames ac ante ipsum primis in
faucibus. Suspendisse potenti. Duis vitae nibh mauris. Duis nec sem sit
amet ante dictum mattis. Suspendisse diam velit, maximus eget commodo
at, faucibus et nisi. Ut a pellentesque eros, sit amet suscipit eros.
Nunc tincidunt quis risus suscipit vestibulum. Quisque eu fringilla
<<<<<<< HEAD:inst/rmarkdown/templates/z-article/skeleton/skeleton.tex
massa. \clearpage
\end{abstract}
=======
massa.
\end{abstract}
}



\begin{document}

\maketitle

%----------------------------------------------------------------------------------------




>>>>>>> 68f0abee1e0db71c1f585da2585915c8eb49dbf6:inst/rmarkdown/templates/zarticle/skeleton/skeleton.tex



\hypertarget{introduction}{%
\section{Introduction}\label{introduction}}

<<<<<<< HEAD:inst/rmarkdown/templates/z-article/skeleton/skeleton.tex
\lettrine[nindent=0em,lines=3]{T} his is my attempt at a full-fledged
RMarkdown article template for use in my package \texttt{zealot}. I have
set some defaults but have made sure to include a lot of options so that
anyone who uses this template from the output format can customize it in
various ways. Really, this whole endeavor is me learning how to
construct a template for myself. Really really, I am procrastinating.
Here is an in-text citation \autocite{campbell1960}. Now I am going to
insert a thing that \textcite{campbell1960} said here just to check.
=======
This is my attempt at a full-fledged RMarkdown article template for use
in my package \texttt{zealot}. I have set some defaults but have made
sure to include a lot of options so that anyone who uses this template
from the output format can customize it in various ways. Really, this
whole endeavor is me learning how to construct a template for myself.
Really really, I am procrastinating. Here is an in-text citation
(Campbell et al. 1960). Now I am going to insert a thing that Campbell
et al. (1960) said here just to check.
>>>>>>> 68f0abee1e0db71c1f585da2585915c8eb49dbf6:inst/rmarkdown/templates/zarticle/skeleton/skeleton.tex

orem ipsum dolor sit amet, consectetur adipiscing elit, sed do eiusmod
tempor incididunt ut labore et dolore magna aliqua. Sed euismod nisi
porta lorem. Nibh venenatis cras sed felis eget velit aliquet sagittis
id. Ultrices tincidunt arcu non sodales. Dictum varius duis at
consectetur. Dictum fusce ut placerat orci. Morbi tristique senectus et
netus et malesuada fames. Arcu cursus vitae congue mauris rhoncus aenean
vel elit. Semper quis lectus nulla at volutpat diam ut venenatis. Nulla
pellentesque dignissim enim sit amet venenatis urna cursus eget.
Suspendisse sed nisi lacus sed viverra tellus in hac habitasse.

Urna nunc id cursus metus. Sit amet massa vitae tortor condimentum
lacinia quis vel eros. Adipiscing commodo elit at imperdiet. Magna etiam
tempor orci eu lobortis elementum nibh tellus molestie. Faucibus vitae
aliquet nec ullamcorper sit. Egestas dui id ornare arcu odio ut sem
nulla pharetra. Velit scelerisque in dictum non consectetur. Tincidunt
tortor aliquam nulla facilisi cras. Augue neque gravida in fermentum et
sollicitudin. Volutpat consequat mauris nunc congue nisi vitae. Quam
pellentesque nec nam aliquam sem et. Justo nec ultrices dui sapien. Mi
quis hendrerit dolor magna. Sagittis eu volutpat odio facilisis mauris
sit. Viverra ipsum nunc aliquet bibendum enim facilisis gravida neque
convallis. Diam quam nulla porttitor massa id neque aliquam.

\hypertarget{methods}{%
\section{Methods}\label{methods}}

\hypertarget{results}{%
\section{Results}\label{results}}

\hypertarget{subsection}{%
\subsection{Subsection}\label{subsection}}

\hypertarget{discussion}{%
\section{Discussion}\label{discussion}}

\hypertarget{subsection-1}{%
\subsection{Subsection}\label{subsection-1}}

\hypertarget{another-subsection}{%
\subsection{Another Subsection}\label{another-subsection}}

<<<<<<< HEAD:inst/rmarkdown/templates/z-article/skeleton/skeleton.tex

\printbibliography
=======
\hypertarget{references}{%
\section*{References}\label{references}}
\addcontentsline{toc}{section}{References}

\hypertarget{refs}{}
\begin{CSLReferences}{1}{0}
\leavevmode\vadjust pre{\hypertarget{ref-campbell1960}{}}%
Campbell, Angus, Philip E. Converse, Warren E. Miller, and Donald E.
Stokes. 1960. \emph{The {American} Voter}. {New York}: {Wiley}.

\end{CSLReferences}
>>>>>>> 68f0abee1e0db71c1f585da2585915c8eb49dbf6:inst/rmarkdown/templates/zarticle/skeleton/skeleton.tex

\end{document}
